\documentclass[11pt]{article}
\usepackage[margin=1in]{geometry}
\usepackage{booktabs}
\usepackage{array}
\usepackage{amssymb}
\usepackage{xcolor}
\usepackage{enumitem}
\usepackage{pifont}
\usepackage{hyperref}
\hypersetup{colorlinks=true, linkcolor=blue, urlcolor=blue}
\newcommand{\code}[1]{\texttt{#1}}
\newcommand{\cmark}{\textcolor{green!55!black}{\ding{51}}}
\newcommand{\xmark}{\textcolor{red!75!black}{\ding{55}}}
\setlength{\parskip}{4pt}
\setlength{\parindent}{0pt}

\title{\textbf{PCCR Parallel Run --- Detailed Report}\\[4pt]
\large A from-scratch explanation of the architecture, the bugs we found,\\
every fix we tried, the full results (accuracy, latency, cost, tokens),\\
and whether the memory stores actually help}
\author{}
\date{2026-06-28 \quad$\cdot$\quad OfficeBench (152 held-out subtasks) \quad$\cdot$\quad Backbone: Cerebras gpt-oss-120b}

\begin{document}
\maketitle

% ======================================================================
\section*{0.\quad TL;DR (read this first)}
\begin{itemize}[leftmargin=1.4em,itemsep=2pt]
\item We built a system (PCCR) = a central memory \emph{router} that decides, per task, \emph{which} of
6 memory stores to consult, instead of always retrieving everything.
\item We compare it against \textbf{retrieve-all} (the LegoMem-style ``always inject the most similar past
episode'') and \textbf{no-memory} (a plain agent).
\item \textbf{Final result: our system $0.421$ (64/152) vs retrieve-all $0.467$ (71/152) --- we are $4.6$
points behind on accuracy, and we use 2$\times$ the tokens (1.5M vs 734K).}
\item We \textbf{win on the hardest tasks ($+3.5$ points)} and on \textbf{document-processing} ($+12$
points), but \textbf{lose on the easy tiers} and especially on \textbf{multi-application} tasks, where our
memory injection falls below even the no-memory agent.
\item The journey fixed \textbf{3 real bugs} (a thread-safety bug, a memory-index pollution bug, and an
API-key-collision bug) that lifted us from $0.368 \to 0.421$, but it did not get us past retrieve-all.
\end{itemize}

% ======================================================================
\section{The architecture, explained from scratch}

\subsection{The problem}
An agent solving an office task (e.g.\ ``extract data from a PDF and email it'') can be helped by
\emph{memory} of past tasks. The naive approach (retrieve-all) injects a similar past episode into the
prompt \emph{on every step}. That helps on easy tasks but (a)~costs a lot of tokens and (b)~can distract
the model on hard tasks. Our thesis: \textbf{route} --- consult a memory store only when it is actually
worth it for \emph{this} task.

\subsection{The six memory stores}
\begin{center}
\begin{tabular}{@{}llp{6.0cm}@{}}
\toprule
Store & What it holds & Everyday analogy \\
\midrule
Procedural memory (PM)  & reusable ``how-to'' rules per task type & a checklist / SOP \\
Working memory (WM)     & the current task's scratchpad (erased at task end) & a sticky note \\
Short-term memory (STM) & cached plans of recently-seen tasks (a hot cache) & ``I just did this one'' \\
Episodic memory (EM)    & full traces of past \emph{successful} tasks, searched by similarity & ``how I did a similar task'' \\
Semantic memory (SM)    & extracted facts & general knowledge \\
Entity memory (ENT)     & the user's durable profile & ``who the user is'' \\
\bottomrule
\end{tabular}
\end{center}
For OfficeBench we have \textbf{62 successful procedures} (built by running the 148 training tasks with no
memory and keeping the successes). These fill the four memory stores above.

\subsection{The router (the core contribution)}
On each task, in the retrieval phase, the router runs a \textbf{cascade}:
\begin{enumerate}[leftmargin=1.6em,itemsep=1pt]
\item \textbf{Procedural and working memory are always read} (cheap, always relevant).
\item \textbf{Short-term memory is checked first.} If the current task is a near-duplicate of a cached
one $\to$ \textbf{short-circuit} (use the cached plan, skip the expensive stores).
\item \textbf{On a miss}, the router \textbf{gates} the three optional stores (episodic, semantic,
entity): consult a store only if it clears a threshold. The original design used a hand-set
``usefulness/cost'' ratio per (task-type, store) --- 15 guessed numbers.
\end{enumerate}

\subsection{What we changed the gate to (the ``confidence gate'' --- our main novelty)}
The 15 hand-set numbers do not generalize. We replace them with \emph{one} rule:
\begin{quote}
consult a store \textbf{only if this query's similarity to that store's content $\ge$ a single threshold $\theta$}.
\end{quote}
So instead of guessing ``episodic memory is worth 0.55 for lookup tasks,'' we \emph{measure} --- if the
actual retrieval is similar enough to the current task, use it. We run with $\theta = 0.45$. (Honest
caveat: $\theta$ is still one hand-chosen number; the principled version is to \emph{learn} it from
outcomes --- implemented but not the headline here.)

\subsection{The other mechanisms we added (to chase accuracy)}
\begin{center}
\begin{tabular}{@{}p{3.6cm}p{8.4cm}@{}}
\toprule
Mechanism & Purpose \\
\midrule
Verified procedure replay & on a high-confidence episodic hit (similarity $\ge 0.75$), adapt the cached
action sequence in \textbf{1 LLM call} and execute it without per-step reasoning; verify each action by
the tool's own success/error signal; \textbf{abort and recover} if anything fails (never blind). \\
Multi-output completeness & replay keeps \emph{all} output-producing steps across apps. \\
Curated procedural rules & one LLM call distils each task type's episodes into a crisp actionable rule + cautions (cached). \\
Output convention & ``write under /testbed/data, exact filename, produce all outputs.'' \\
Completion gate & block the ``finish'' action until every output file the task names exists. \\
Self-verification & one reflection call before finishing (``is everything correct/complete?''). \\
Malformed-action recovery & when the agent emits an invalid action, inject the valid action names for that app. \\
Raised step budget & the hardest tier's step cap raised $30 \to 45$. \\
\bottomrule
\end{tabular}
\end{center}

% ======================================================================
\section{The PARALLEL architecture (the ``parallel thing'')}
The reference system runs sequentially. We added concurrency in \textbf{two dimensions} (only the first is
used on OfficeBench):
\begin{itemize}[leftmargin=1.4em,itemsep=1pt]
\item \textbf{Dimension 1 --- parallel TASKS:} run $N$ tasks at once. A task queue (semaphore of size $N$)
pulls from one pool of all 152; $N$ workers run concurrently.
\item \textbf{Dimension 2 --- parallel SUB-AGENTS} (within one task): a dependency analyzer turns
sub-tasks into a DAG of ``waves.'' \emph{Proven equivalent to sequential on a deterministic harness, but
NOT used on OfficeBench} (its agent is inherently one-action-at-a-time).
\end{itemize}

\subsection{The OfficeBench parallel adapter}
OfficeBench runs each task in a \textbf{Docker container}. The blocking per-task routine (Docker exec +
LLM HTTP) is wrapped so $N$ run concurrently:
\begin{itemize}[leftmargin=1.4em,itemsep=1pt]
\item Run the blocking work in a worker thread so the event loop can dispatch the other $N{-}1$ tasks
(otherwise they serialize and there is no speedup).
\item One Docker container per worker (\code{ob-test-0..3}) --- each task gets an isolated \code{/testbed}.
(All tasks use the same paths, so they \emph{cannot} share one container --- they would overwrite each
other's files and corrupt the evaluation.)
\item Longest-task-first ordering so the long hard tasks do not tail-end and stall a worker.
\item One shared, read-only copy of the memory (the 62 procedures) across all threads.
\end{itemize}
We measured a \textbf{$\sim$2.35--3.3$\times$ speedup at $N{=}4$} (ceiling 4$\times$; the gap is tail
effects, per-task Docker setup, a single shared Docker daemon, and rate-limit stalls).

% ======================================================================
\section{The first run, and what we found}
The first full $N{=}4$ run scored \textbf{0.368} (56/152) vs retrieve-all 0.467 --- 10 points behind. We
categorized all 96 failures:
\begin{center}
\begin{tabular}{@{}rl@{}}
\toprule
\# & Failure mode \\
\midrule
29 & malformed-loops --- the model emits a WRONG action name and loops to the step cap \\
41 & wrong-content --- produced an output, but content incomplete/wrong (mostly hard multi-app) \\
21 & missing-output --- quit before creating a required file \\
5  & wrong-cell --- wrote to the wrong spreadsheet row/column \\
\bottomrule
\end{tabular}
\end{center}
This first run turned out to be confounded by three bugs (next section); once found, the ``this feature
hurts'' conclusions we had drawn were mostly artifacts of these bugs.

% ======================================================================
\section{The three real bugs we found and fixed}
\begin{enumerate}[leftmargin=1.5em,itemsep=3pt]
\item \textbf{A thread-safety bug in the benchmark's timeout.} OfficeBench wraps every action in a
timeout that uses an OS signal --- and OS signals only work in the main thread. Our parallel adapter runs
tasks in worker threads, so the timeout threw on \emph{every} action $\to$ reported as ``Malformed
action!'' $\to$ 89\% malformed. \textbf{Fix:} make the timeout main-thread-aware. Verified $89\% \to 0\%$.
\item \textbf{Memory-index pollution (the biggest bug).} The similarity index used a shared, persistent
folder that \emph{loaded existing vectors and ADDED 62 more on every build} ($62\to124\to186\to\ldots$).
Crucially, \textbf{retrieve-all used a fresh clean index, but our system used the polluted one} --- an
unfair handicap. \textbf{Fix:} a fresh isolated index every time $\to$ exactly 62 clean procedures. This
alone lifted \textbf{the easy tier from 0.468 to 0.617}.
\item \textbf{API-key collision.} Every task's LLM client started at key \#1, so $N$ concurrent tasks
marched in lockstep and hit the same key every step $\to$ rate-limit storms. \textbf{Fix:} a global
round-robin so concurrent calls get different keys. Cut the rate limits sharply and improved stability.
\end{enumerate}
\textbf{Important:} the malformed actions \emph{themselves} are gpt-oss nondeterminism, NOT our code ---
proven by running the same task, same code, sequentially three times: \textbf{80\% / 0\% / 62\%
malformed}. The provider flakes on hard tasks. Our malformed-action recovery mitigates this but cannot
eliminate it.

% ======================================================================
\section{Final results (full)}
\subsection{Accuracy}
\begin{center}
\begin{tabular}{@{}lccc@{}}
\toprule
 & our system & retrieve-all & no-memory \\
\midrule
ALL (152) & \textbf{0.421} & \textbf{0.467} & 0.414 \\
L1 (47)   & 0.617 & 0.723 & 0.553 \\
L2 (48)   & 0.500 & 0.583 & 0.500 \\
L3 (57)   & \textbf{0.193} & 0.158 & 0.228 \\
\bottomrule
\end{tabular}
\end{center}
Win/Loss/Tie vs retrieve-all: 11 / 18 / 123. (We beat retrieve-all on the hardest tier.)

\subsection{By task type (the most revealing cut)}
\begin{center}
\begin{tabular}{@{}lccc@{}}
\toprule
Task type & our system & retrieve-all & no-memory \\
\midrule
document-processing & \textbf{0.720} & 0.600 & 0.480 \quad(\textbf{we WIN}) \\
lookup              & 0.750 & 0.812 & 0.625 \\
single-action       & 0.565 & 0.609 & 0.478 \\
data-compute        & 0.700 & 0.800 & 0.700 \\
multi-application    & \textbf{0.179} & 0.269 & 0.295 \quad(\textbf{we LOSE --- below no-memory}) \\
\bottomrule
\end{tabular}
\end{center}
\textbf{Story:} memory clearly helps document-processing, lookup, single-action and data-compute. But on
\textbf{multi-application tasks (78 of 152 --- half the benchmark)}, our injection drops below no-memory ---
the heavy injection + the gpt-oss malformed-loops actively hurt the hardest multi-app tasks. This single
task type is why we lose overall.

\subsection{Latency / throughput}
\begin{center}
\begin{tabular}{@{}ll@{}}
\toprule
avg wall-clock / task & 27.2 s \\
avg compute / task (rate-limit time removed) & 24.8 s \\
total rate-limit hits & 506 (down sharply after the key-rotation fix) \\
parallel speedup @ $N{=}4$ & $\sim$2.35--3.3$\times$ (ceiling 4$\times$) \\
\bottomrule
\end{tabular}
\end{center}
Docker note: the Docker VM crashed 3$\times$ under 4-container load; we resumed from checkpoints each time.

\subsection{Cost (tokens) --- the weak point}
\begin{center}
\begin{tabular}{@{}lr@{}}
\toprule
our system total injected tokens & 1{,}508{,}634 \\
retrieve-all & 734{,}295 \quad($\to$ we use $\sim$2$\times$ MORE) \\
avg LLM calls / task & 17.5 (inflated by rate-limit retries) \\
avg agent steps / task & 14.0 \\
\bottomrule
\end{tabular}
\end{center}

% ======================================================================
\section{Why so many tokens, and what to do about it}
The memory block is re-injected on every agent step. Per retrieval it breaks down as (approx tokens):
episodic memory (past plans + actions) 255; curated procedural rule 183; output convention 147; semantic
facts 59; step-budget hint 12. So \textbf{the curated rule + convention add $\sim$330 tokens/step that
retrieve-all does not have}. On top of that, the completion gate + self-verification add extra LLM calls
that re-inject the whole block; and malformed-loops run many steps, each re-injecting $\to$ huge inflation
on the multi-app/hard tasks that fail anyway.

\textbf{What to do better:}
\begin{enumerate}[leftmargin=1.5em,itemsep=1pt]
\item Inject the static parts once, not every step (the procedural rule + convention do not change within a task).
\item Skip the episodic text-hint when replay fires (replay already executes the procedure).
\item Drop the curated rule + convention for a cost-optimal arm: confidence-gated episodic memory alone is
\emph{cheaper} than retrieve-all's always-on episodic memory.
\item Cap re-injection on malformed-loops (runaway cost is on tasks that fail anyway).
\end{enumerate}

% ======================================================================
\section{Are the memory stores actually helping?}
\begin{itemize}[leftmargin=1.4em,itemsep=2pt]
\item \textbf{Procedural memory: YES, broadly used and helpful.} Fired on \textbf{152/152}; lifts
document-processing $+24$ points over no-memory. \emph{How written:} one offline LLM call distils each
task type's 62 episodes into a 2--4 sentence rule with cautions.
\item \textbf{Episodic memory: YES on easy/similar tasks, NEUTRAL/negative on multi-app.} Consulted on
\textbf{121/152} (gate declined 31 low-similarity tasks). Drives lookup/single-action/data-compute. On
multi-app the episode does not transfer and can mislead (stale row index).
\item \textbf{Short-term memory: NOT helping accuracy here --- by design.} Only \textbf{8/152}
short-circuits (single-pass test has almost no recurrences). It is a latency/cost mechanism for repeat
workloads.
\item \textbf{Verified replay: fired on 29/152} high-confidence tasks --- cuts LLM calls and won 4 of the
11 head-to-head rescues.
\item \textbf{Semantic \& entity memory:} the gate mostly suppresses them (noise on procedural office
tasks); entity memory is empty on OfficeBench.
\end{itemize}

% ======================================================================
\section{Three real task lifecycles (the latest 64/152 architecture)}
Each shows the full pipeline: ingest $\to$ retrieve (gate) $\to$ execute $\to$ finish-check $\to$ outcome.

\subsection*{Lifecycle A --- Replay win with verified abort-and-recover $\cdot$ task 3-60/0 (multi-app, hardest tier) \cmark}
\textbf{Task:} ``Find all students taking CS161, put name + student ID into an output file.''
\begin{enumerate}[leftmargin=1.5em,itemsep=1pt]
\item \textbf{Ingest:} task text $\to$ working memory; classified as a multi-application task.
\item \textbf{Retrieve / gate:} short-term miss. The gate measures episodic similarity $= 0.832 \ge
\theta(0.45) \to$ consult episodic memory (semantic too; entity empty). The procedural rule is injected.
\item \textbf{Replay:} similarity $0.832 \ge 0.75 \to$ replay fires; one adapt LLM call rewrites the best
cached procedure (11 actions) for this task.
\item \textbf{Verified execution:} runs 5 actions, then the cached path \code{CS161\_roster.csv} does not
exist (real file \code{CS161\_Class\_Roster.xlsx}) $\to$ tool error $\to$ \textbf{replay ABORTS} (never blind).
\item \textbf{Recovery:} the agent re-reads the data folder, finds the real file, filters CS161, writes output.
\item \textbf{Finish-check:} the completion gate confirms the output exists $\to$ allowed. \textbf{SUCCESS}
(45 steps, 52 calls, 41{,}960 tokens).
\end{enumerate}
\emph{Lesson:} the whole thesis in one task --- memory gives a head start (replay), tool-signal
verification catches the mismatch, the agent recovers. ``Memory as an executable, monitored procedure.''

\subsection*{Lifecycle B --- Short-term cache short-circuit $\to$ cheap win $\cdot$ task 3-59/0 (multi-app, hardest tier) \cmark}
\textbf{Task:} ``Analyze students' grade data, generate a teaching report in teaching.docx.''
\begin{enumerate}[leftmargin=1.5em,itemsep=1pt]
\item \textbf{Retrieve / gate:} short-term memory checked first $\to$ \textbf{HIT} (near-duplicate). The
cascade \textbf{short-circuits}: the expensive stores are NOT searched $\to$ only procedural + working +
short-term are used; the cached plan is injected directly.
\item \textbf{Replay:} independently finds a near-identical episode (similarity $0.945$), adapts, executes
4, aborts on a filename mismatch, the agent recovers and creates \code{teaching.docx}.
\item \textbf{Finish-check:} confirms the file $\to$ finishes. \textbf{SUCCESS} in 19 steps and only
\textbf{3{,}872 tokens} ($\sim$10$\times$ cheaper than A).
\end{enumerate}
\emph{Lesson:} short-term memory is the cost lever --- on a recurrence it skips the whole similarity
search. Only 8 fired here, but on a recurring workload it would dominate.

\subsection*{Lifecycle C --- Multi-application failure (the task type that sinks us) $\cdot$ task 3-40/0 \xmark}
\textbf{Task:} ``Collect car trading records into car\_records.xlsx, AND schedule a meeting'' (2+ outputs).
\begin{enumerate}[leftmargin=1.5em,itemsep=1pt]
\item \textbf{Retrieve / gate:} short-term miss. Episodic similarity $= 0.371 < \theta(0.45) \to$
\textbf{episodic memory GATED OUT} (no good episode); replay does not fire. Only procedural + working
memory are used --- the gate correctly declining a weak memory.
\item \textbf{Execution:} the model reasons from scratch; emits \code{shell/list\_directory} --- but the
valid shell action is \code{command} $\to$ ``Malformed action!''. The recovery injects valid names, but
gpt-oss keeps emitting the invalid name (stuck loop --- provider nondeterminism), then gives up.
\item \textbf{FAILURE} in 5 steps (the required \code{car\_records.xlsx} was never created).
\end{enumerate}
\emph{Lesson:} exactly why multi-app is below no-memory (0.179 vs 0.295). On hard tasks with no good
episode (correctly gated out), the agent is on its own, and gpt-oss's malformed-loop flakiness (not our
code) tanks it.

\begin{center}
\begin{tabular}{@{}lcccc@{}}
\toprule
 & Episodic sim & Gate decision & Mechanism & Tokens & Outcome \\
\midrule
A (3-60) & 0.832 & consult episodic+semantic & replay$\to$abort$\to$recover & 41{,}960 & \cmark \\
B (3-59) & 0.945 & short-term short-circuit & cached plan + replay & 3{,}872 & \cmark \\
C (3-40) & 0.371 & episodic gated out & step-by-step (no memory) & 930 & \xmark \\
\bottomrule
\end{tabular}
\end{center}

% ======================================================================
\section{Hard evidence: which tasks the memory actually rescued}
\emph{(The ``no-memory'' column proves the memory --- not the base agent --- made the difference.)}

\subsection{The 11 tasks where our system PASSED but retrieve-all FAILED}
\begin{center}
\begin{tabular}{@{}llllp{4.6cm}@{}}
\toprule
Task & Tier & Task type & What won it & no-memory \\
\midrule
1-2/1  & L1 & single-action  & episodic + procedural & \textbf{FAIL} \\
1-8/2  & L1 & single-action  & \textbf{replay} + episodic + procedural & \textbf{FAIL} \\
2-1/3  & L2 & document-proc  & episodic + procedural & \textbf{FAIL} \\
2-20/2 & L2 & multi-app      & \textbf{replay} + episodic + procedural & \textbf{FAIL} \\
3-10/0 & L3 & multi-app      & episodic + procedural & \textbf{FAIL} \\
3-60/0 & L3 & multi-app      & \textbf{replay} + episodic + procedural & \textbf{FAIL} \\
3-85/0 & L3 & document-proc  & \textbf{procedural rule alone} & \textbf{FAIL} \\
3-95/0 & L3 & document-proc  & \textbf{replay} + episodic + procedural & \textbf{FAIL} \\
2-32/1 & L2 & data-compute   & procedural & pass \\
3-8/3  & L3 & multi-app      & procedural & pass \\
3-9/1  & L3 & multi-app      & procedural & pass \\
\bottomrule
\end{tabular}
\end{center}
\textbf{Read this:} 8 of 11 are ``no-memory FAIL'' $\to$ genuine memory rescues retrieve-all missed.
Replay won 4 (1-8, 2-20, 3-60, 3-95), episodic+procedural won 3, and \textbf{3-85 was won by the
procedural rule alone} (episodic gated out, no-memory failed).

\subsection{Procedural memory --- the always-on safety net}
Fires 152/152. Clearest proof: \textbf{document-processing: our system 18/25 vs no-memory 12/25 vs
retrieve-all 15/25} ($+6$ over no-memory, $+3$ over retrieve-all). \textbf{Task 3-85/0} is the smoking
gun: episodic gated out, no-memory failed, the procedural rule alone got it to PASS.

\subsection{Short-term memory --- the cost lever, proven on recurrences}
\begin{center}
\begin{tabular}{@{}lllll@{}}
\toprule
Task & Task type & Tokens & Result & Note \\
\midrule
2-22/0 & document-proc & \textbf{484}  & \cmark & $\sim$3$\times$ cheaper than normal \\
2-39/0 & multi-app     & \textbf{466}  & \cmark & \\
2-21/3 & multi-app     & 2{,}178       & \cmark & \textbf{recurrence of 2-21/0} --- cache caught the duplicate \\
2-21/0 & multi-app     & 2{,}662       & \cmark & ``extract Bill's courses\ldots'' \\
3-59/0 & multi-app     & 3{,}872       & \cmark & (Lifecycle B) \\
3-19/0, 3-20/0, 3-24/0 & multi-app & $\sim$5--7k & \xmark & hard PDF tasks; cache hit but task still hard \\
\bottomrule
\end{tabular}
\end{center}
\textbf{Read this:} the cache delivered \textbf{5/8 successes at a fraction of the token cost} (down to
$\sim$466 tokens). 2-21/0 and 2-21/3 are duplicate phrasings of the same task --- the cache recognized the
recurrence and reused the plan, exactly its purpose. On a recurring production workload the cache would
short-circuit the majority of tasks --- where the big cost/latency win lives.

% ======================================================================
\section{What we can derive (the honest state of the project)}
\begin{enumerate}[leftmargin=1.5em,itemsep=3pt]
\item \textbf{The architecture is sound and does the right things.} The gate consults memory when
confident and declines when not; replay executes-and-verifies; procedural memory is a reliable net;
short-term memory short-circuits recurrences. We have named, reproducible task-level evidence for each.
\item \textbf{Memory genuinely helps where it should:} document-processing ($+6$ vs no-memory), lookup,
single-action, data-compute, and 8 specific rescues retrieve-all missed. The memory-index fix lifted us
$+5.3$ points overall and $+15$ on the easy tier.
\item \textbf{Two honest weaknesses keep us behind retrieve-all (0.421 vs 0.467):}
  \begin{itemize}[leftmargin=1.2em,itemsep=0pt]
  \item \emph{Accuracy:} the multi-application task type (78 tasks) drops below no-memory (0.179 vs 0.295)
  --- hard tasks with no good episode + gpt-oss malformed-loop nondeterminism (proven not our code).
  \item \emph{Cost:} we inject 2$\times$ retrieve-all's tokens (1.5M vs 734K) due to the curated rule +
  convention + the gates + malformed-loop re-injection.
  \end{itemize}
\item \textbf{The clear path to ``better accuracy, lower cost'':}
  \begin{itemize}[leftmargin=1.2em,itemsep=0pt]
  \item \emph{Cost:} drop the curated rule + convention to a lean arm (confidence-gated episodic memory
  only $\to$ cheaper than retrieve-all), inject static parts once, skip the episodic hint when replay fires.
  \item \emph{Accuracy:} rescue multi-application --- stronger malformed-action retry, do not over-inject
  on low-confidence tasks, lean on the completion gate to force the 2nd/3rd outputs, learn the threshold
  from outcomes. If multi-app merely reaches no-memory parity, the overall jumps $\sim+6$ points $\to$
  neck-and-neck with retrieve-all, at lower cost.
  \end{itemize}
\item \textbf{The most defensible claim today} is not ``beats retrieve-all on everything.'' It is:
\emph{``A routing memory that is cost-aware and verification-safe --- it matches or beats retrieve-all on
structured tasks (document-processing) and on hard tasks (where always-retrieve collapses), rescues tasks
that no-memory and retrieve-all both miss, and turns recurrences nearly free via the short-term cache ---
while the remaining gap is a backbone-stability problem on open-ended multi-application tasks, not a
routing flaw.''}
\end{enumerate}

\vspace{6pt}
\noindent\textbf{Reproduce:} \code{source cerebras.env \&\& python -m officebench\_eval.run\_t1\_parallel
--concurrency 4 --improved} \quad(traces in \code{traces/t1\_par/}, results in \code{t1\_parallel\_progress.json}).

\end{document}
